\chapter{People: invitations and signup links}
\label{ch:people}

\section{How accounts work}

There is no automatic connection between the university's user directory and the chat
server. Accounts are created through \emph{signup links} that the bot issues and that you (or
an admin) send to people.

\begin{itemize}
  \item \textbf{The username comes from the university email address.} The part before the
        \cmd{@} becomes the username: \cmd{anna.beispiel@stud.unibas.ch} becomes
        \cmd{@anna.beispiel:matrix.dmi.unibas.ch}. Nobody can choose a different name, so you
        can always tell who an account belongs to.
  \item \textbf{\cmd{@unibas.ch} and \cmd{@stud.unibas.ch} are the same person.} The same name
        on both domains leads to the same account. If one list uses the student address and
        another the staff address, nobody gets two accounts.
  \item \textbf{A link works once and expires after 14 days.} When it is used, the account is
        created with the password the person chooses, and they are added to the course the
        link was made for.
  \item \textbf{The bot doesn't send emails.} It gives you the links; you send them, e.g.\ by
        mail merge from Outlook. Your own address as sender is also more trustworthy for
        students than an unknown system address.
\end{itemize}

Only university addresses are accepted. Addresses with unusual characters before the
\cmd{@} (such as \cmd{+}) are rejected, because they can't be turned into a username safely.
Guests without a university address need a bot admin (\autoref{sec:guests}).

\section{Invite a whole list: \texttt{!invite}}

\begin{you}
!invite CS101
anna.beispiel@stud.unibas.ch
max.muster@stud.unibas.ch, eva.example@unibas.ch
someone@gmail.com
\end{you}

Put the addresses below the command, one per line or separated by commas or spaces. For
long lists, send just \cmd{!invite CS101} and then the list as a \file{.txt} or \file{.csv}
file, e.g.\ the participant export of the course platform. The bot picks out everything that
looks like an email address, so extra columns don't matter.

\begin{botreply}
Invite to \textbf{CS101}: 1 existing accounts get an invitation, 2 new people get a signup
link (as a CSV file after you confirm), 1 entries are not usable.\\
Not usable: someone@gmail.com\\[4pt]
Nothing has been changed yet. To run this, send \cmd{!confirm k7mq} within 10 minutes. To
drop it, send \cmd{!cancel}.
\end{botreply}

After \cmd{!confirm k7mq}:

\begin{itemize}
  \item People who already have an account get an \textbf{invitation} to the course space.
        They see it in Element and accept it; then they can join all open course rooms.
  \item For new people the bot creates \textbf{signup links} and sends you
        \file{signup-links.csv} with the columns \cmd{email}, \cmd{user_id}, \cmd{link} and
        \cmd{course}, plus a text you can use as the email.
  \item When a new person uses their link, the account is created and they are
        \textbf{added to the course space and its open rooms right away}; they don't have
        to accept anything.
\end{itemize}

\begin{tip}
Use \cmd{--dry-run} to check a list first: you see the counts and the unusable entries
without creating a plan.
\end{tip}

\subsection{Send the links by mail merge}
\label{sec:mailmerge}

Each person must get \emph{their own} link: it creates exactly the account named in it.

\begin{enumerate}
  \item Save \file{signup-links.csv} from the chat.
  \item In Word: \ui{Mailings} \then \ui{Start Mail Merge} \then{}
        \ui{E-mail Messages}, then \ui{Select Recipients} \then \ui{Use an
        Existing List} and choose the CSV file.
  \item Write the message (the bot suggests a text) and insert the field \cmd{link} with
        \ui{Insert Merge Field}.
  \item \ui{Finish \& Merge} \then \ui{Send E-mail Messages}, with \cmd{email}
        as the recipient column.
\end{enumerate}

Other tools work the same way: one message per row, with the \cmd{link} column inserted.
For a single person, copy the link into a normal email.

\begin{warning}[The links are personal]
Don't post the CSV file or a link in a course room or forward it to the wrong person: a link
creates the account named in it for whoever opens it first.
\end{warning}

\section{Single links: \texttt{!signup}}

For latecomers or a single new TA:

\begin{you}
!signup new anna.beispiel@stud.unibas.ch CS101
\end{you}
\begin{botreply}
Signup link for \textbf{@anna.beispiel:matrix.dmi.unibas.ch} (anna.beispiel@stud.unibas.ch):\\
https://matrix.dmi.unibas.ch/\_matrix/maubot/plugin/adminbot/signup?t=\dots\\[4pt]
Text to forward: \dots
\end{botreply}

With a course code, the person is added to that course when they sign up. If they already
have an account, the bot invites them to the course right away instead. Without a course
code (professors only), the link only creates the account.

\begin{center}
\begin{tabularx}{\textwidth}{@{}lL@{}}
\toprule
\textbf{Command} & \textbf{Use it when} \\
\midrule
\cmd{!signup list [course]} & you want to see who hasn't used their link yet (with expiry date) \\
\cmd{!signup renew <email>} & a link got lost or is about to expire: one new link, old ones stop working \\
\cmd{!signup revoke <email>} & a link went to the wrong address: all open links stop working \\
\bottomrule
\end{tabularx}
\end{center}

Instead of the email you can also give the Matrix ID (\cmd{@anna.beispiel:matrix.dmi.unibas.ch}).
You can manage links you created yourself and the links of courses where you are staff.

\section{Guests (admins)}
\label{sec:guests}

Guests without a university address, e.g.\ a guest lecturer, need a bot admin:
\cmd{!signup guest <email> <name> [course]} creates a link for the account
\cmd{@guest.<name>}.

\section{Passwords and lost keys}
\label{sec:passwords}

There is no \enquote{forgot password} email on \server: the server sends no emails. Instead,
everyone, students included, can reset their own password with the bot. You don't have to do
anything for your students except point them here.

\subsection{Three ways to a new password}

\begin{center}
\begin{tabularx}{\textwidth}{@{}lLL@{}}
\toprule
\textbf{Situation} & \textbf{What to do} & \textbf{Sessions} \\
\midrule
Still logged in somewhere (e.g.\ the phone), forgot the password &
  Send \cmd{!password} to the bot. It replies with a link for the own account, valid
  1~hour. & stay logged in \\
Logged out everywhere, has the reset code &
  Open the reset page (\url{https://matrix.dmi.unibas.ch/_matrix/maubot/plugin/adminbot/recover}),
  enter username, reset code and a new password. The page then shows a new code. &
  all logged out \\
Neither &
  A bot admin creates a reset link with \cmd{!signup reset @user} (24~hours). &
  all logged out \\
\bottomrule
\end{tabularx}
\end{center}

\begin{you}
!password
\end{you}
\begin{botreply}
Here is a link to set a new password for \textbf{@anna.beispiel:matrix.dmi.unibas.ch}
(valid 60 minutes, works once):\\
https://matrix.dmi.unibas.ch/\_matrix/maubot/plugin/adminbot/reset?t=\dots\\[4pt]
Your sessions stay logged in. \dots
\end{botreply}

\begin{description}[style=nextline,leftmargin=1.5em,font=\normalfont\bfseries]
  \item[Reset code]
    Shown once on the signup page when the account is created. \cmd{!password code} gives a
    new one (the old one stops working); people with accounts from before version 0.3 get
    their first code this way. Each code works once. It must be kept outside Matrix, in a
    password manager or on paper, because whoever needs it can't read their chats.
  \item[Notices]
    After every reset the bot tells the person in a direct chat. If all their sessions were
    logged out, the message waits until they have logged in again, so they can read it.
  \item[Why professors can't reset students' passwords]
    Whoever creates a reset link for someone else can open it and take over the account.
    That's why only bot admins can do it, and why the self-service ways only ever work for
    the sender's own account.
\end{description}

\subsection{Recovery key: why old messages can be lost}

Element's \emph{recovery key} unlocks the encrypted messages. It is created on the person's
device and never reaches the server or the bot, so \textbf{nobody can restore it}. A new
password doesn't help with this:

\begin{itemize}
  \item \textbf{With the recovery key}, a new login just asks for it, and everything is
        readable as before.
  \item \textbf{Without it}, Element asks to confirm the new login. The way out is
        \ui{Can't confirm?} and then \ui{Reset identity} (the words differ a little between
        Element versions; it needs the account password). Afterwards the person can chat
        normally and is shown as verified again, but older encrypted messages stay
        unreadable. Element then shows a new recovery key: this time, save it.
\end{itemize}

\begin{tip}[Text for your students]
Paste this into your first announcement:\par\smallskip
\enquote{When you log in to Element for the first time, set up \emph{recovery} and keep the
recovery key in a password manager or on paper, together with the reset code from the signup
page. Forgot your password? If you're still logged in somewhere, send \cmd{!password} to
@adminbot. Otherwise use your reset code on the reset page. Lost the code? \cmd{!password
code} gives you a new one.}
\end{tip}

\section{Compare with the official list: \texttt{!roster sync}}

Over the first weeks, people drop out and others join late. \cmd{!roster sync} compares a
current participant list with the people in the course space:

\begin{you}
!roster sync CS101
\end{you}

Then send the list as a file (or put it below the command, as with \cmd{!invite}).

\begin{botreply}
Roster \textbf{CS101}: 120 on the list, 112 already in the course, 8 missing, 5 in the course
but not on the list.\\
Missing: \dots\ Invite them with \cmd{!invite CS101} and the same list.\\
Not on the list: \dots\\
To remove them, send the same command with \cmd{--remove}. Nothing was changed.
\end{botreply}

Staff (owner and TAs) are never counted or removed. With \cmd{--remove}, the people not on
the list are removed from the course space and all its rooms, after you confirm the plan.
